\scanpage{167}

\uline{Ex} Si $X$ \deleted{\unsure{est} \illegible{1} \unsure{contient} \unsure{le} \illegible{2} irréductibles} \inserted{est \uline{réduit}, alors schématiquement dense = dense}, pseudo-morphismes $=$ applications rationnelles

Pour tout ouvert $U$ de $X$, soit $F(U)$ l'ens des pseudo-$S$-morphismes de $U$ dans $Y$ ($S$-schém. séparé). Application \deleted{de}restriction $F(U)\unsure{\mathord{-}|\longrightarrow}F(\unsure{V|U})$ si $U\supset V$. Cela \deleted{est} définit un \uline{faisceau} [grâce à la relation $\mathrm{Dom}(f_i\mid U_{ij})=\mathrm{Dom}(f_i)\cap U_{ij}$], appelé \uline{faisceau des} \deleted{\unsure{quasi-morph.}} \uline{pseudo-$S$-morphismes} de $X$ dans $Y$.

En particulier, on obtient ainsi le faisceau des pseudo-sections de $\underline O_X$ (\deleted{\unsure{f}}faisant $X=S$, $Y=X[t]$) qui est un faisceau d'anneaux \uline{\unsure{contenant|$\longleftarrow$}} $\underline O_X$, en particulier un faisceau d'algèbres.

\deleted{\uline{Prop.} 0.1 Ce faisceau est quasi-cohérent. On peut supposer $X=\mathrm{Spec}(A)$.}

% archive sheet number: 166
